\chapter{Weak Topologies}
\vspace{-25pt}
Let $X$ be a LCS. For any $x \in X$, recall that the \textit{evaluation map} is $\hat{\cdot}:X^\ast \rightarrow \bfR$ given by $\hat{x}(f) := f(x)$. We introduce the following notation for convenience.
\mbox{}
\begin{itemize}[itemsep=1pt,topsep=3pt]
	\item For $f \in X^\ast$, define $\left<\cdot, f \right>:X \rightarrow \bfF$ by $x \mapsto f(x)$.
	\item For $x \in X$, define $\left<x,\cdot \right>:X^\ast \rightarrow \bfF$ by $f \mapsto \hat{x}(f) = f(x)$.
\end{itemize}
We call $\left<\cdot,\cdot \right>$ the \textit{canonical evaluation map}.

\section{Basic Definitions and Examples}

\begin{definition}
	Let $X$ be a LCS.
	\mbox{}
	\begin{enumerate}[label = (\alph*),itemsep=1pt,topsep=3pt]
		\item The \textui{weak topology} on $X$, denoted $\sigma (X,X^\ast)$ (or $wk$), is the topology generated by the family of seminorms $\{p_f:X \rightarrow \bfF : f \in X^\ast\}$, where $p_f(x) := |\left<x,f \right>|$.
		\item \textui{weak-* topology} on $X^\ast$, denote $\sigma(X^\ast,X)$ (or $wk*$), is the topology generated by the family of seminorms $\{p_x:X^\ast \rightarrow \bfF :x \in X\}$, where $p_x(f) := |\left<x,f \right>|$.
	\end{enumerate}
\end{definition}

\begin{remark}
	Note that $\sigma(X,X^\ast)$ has $\{p_f^{-1}(U) : f \in X^\ast, U \subseteq \bfF \text{ open }\}$ as a subbase. As such, given $f_1,\ldots ,f_k \in X^\ast$ and $U_1,\ldots U_k$ open in $\bfF$, an arbitrary basis element is of the form $\bigcap_{i = 1}^k f_k^{-1}(U_k)$.
\end{remark}



\begin{proposition}
	Let $X$ be a LCS with the weak topology. Then $U$ is open in $X$ if and only if for every $x_0 \in U$ there exists $\epsilon > 0$ and $f_1,\ldots ,f_n \in X^\ast$ such that:
	\[
		\bigcap_{k = 1}^n \{x \in X : |\left<x-x_0,f_k \right>| < \epsilon\} \subseteq U.
	\] 
\end{proposition}
\iffalse
	\begin{proof}
		Let $U$ be an open set of $(X,wk)$. Let $x_0 \in U$. There exists $f_1,\ldots ,f_k \in X^\ast$ and $V_1,\ldots ,V_k$ open in $\bfF$ such that:
		\[
			x_0 \in \bigcap_{i=1}^k p_{f_i}^{-1}(V_i) \subseteq U.
		\]
		Since $V_i$ is open in $\bfF$, for each $i=1,\ldots ,k$ there exists an $\epsilon_i > 0$ such that:
		\[
			p_{f_i}^{-1} \left( B(p_{f_i},\epsilon_i) \right) \subseteq p_{f_i}^{-1}(V_i) 
		\] 
		We can rewrite the left set as:
		\[
			p_{f_i}^{-1}\left( B(p_{f_i},\epsilon_i) \right) = \{x \in X : |p_{f_i}(x) - p_{f_i}(x_0)| < \epsilon_i\}. 
		\] 
		Now define $\epsilon := \min \{\epsilon_1,\ldots ,e_k\}$. Since $p_{f_i}$ is a seminorm for all $i=1,\ldots ,k$, the reverse-triangle inequality gives:
		\[
			|p_{f_i}(x) - p_{f_i}(x_0)| \leq p_{f_i}(x - x_0) < \epsilon < \epsilon_i.
		\] 
	\end{proof}
	Answer briefly and directly. Give the conclusion first. Do not add background, analogies, or extra explanation unless I explicitly ask. For math questions: avoid restating the question, avoid filler, avoid hedging, and avoid motivational language. After answering a question, do not repeat yourself under any circumstances. After answering a question, do not summarize your explanation under any circumstances.
\fi
